% Pursuing Stacks (cotes 134-2 to 134-8): archive cote 134-6, site page 62 (PDF page 63), typescript page 380.
% typescript. Draft compared: ps5.tex:4506-4572.
% AI-transcribed by Claude Opus 5.5 (claude-opus-5-5, adaptive thinking, effort high), Message Batches API, batch ps3 (msgbatch_01XdFmuVVBAcgc7EFxdtCqsj), 2026-11-12; not yet checked by eye.
\scanpage{134-6/62}
\mspage{380}

\inserted{abelianization of} homotopy types modelized by a var\xout{j}iable object \inserted{$X$} in $A^\wedge$, and b) we want that
the chain complex in (Ab) expressing abelianization of X, should be
expressible in terms of the "tautological abelianization" $\mathrm{Wh}_A(X) = \underline{Z}^{(X)}$
of X , \inserted{itself} by a formula \inserted{moreover} which should make sense functorially with respect
to an arbitrary abelian presheaf \unsure{i.e.} \inserted{an} object F in $A^\wedge_{ab}$ .

The main fact I have in view here is that whenever the chain comple\unsure{x}
L. in Addinf(B) is endowed with an augmentation
\[ (10) \qquad L_{\bullet} \longrightarrow \underline{Z}_B \]
turning it into a resolution of $\underline{Z}_B$ i.e. into a (quasi special) \uline{integra-}%
\uline{tor} for B, then ipso facto L. is an abelianizator for A, the commuta-%
tion isomorphism \hand{$\lambda$} being canonically defined by the augmentation (10).

Some comments, before proc\inserted{e}eding to a proof. Presumably, the
converse of our statement holds too - namely that the natural functor
we'll get from quasi-special integrators for B to abelianizators \inserted{$(L_{\bullet},\lambda)$} for A
is an equivalence \inserted{(even an isomorphism !)} between the relevant categories. I do'nt feel like
pursuing this - \xout{more} the more relevant fact here, whether or not a con-%
verse as contemplated holds, is that we can pin down at any rate a
special class of \xout{inte} abelianizators for A, namely those which come from
\inserted{(quasi special)} integrators for B, and these abelianizators are defined up to \inserted{chain} homotopism
in Addinf(B). In this sense, \uline{we get an existence and unicity statement}
\inserted{\uline{for abelianiz-ators in $A$},} as strong as we possibly could hope for. In practical terms, it would
seem, \uline{an abelianizator for A will be no more no less than just a\xout{n}} \inserted{(quasi-special)} \uline{inte-}%
\uline{grator for B}, namely \uline{a projective resolution of $\underline{Z}_B$ in $B^\wedge_{ab}$} , \xout{having}
\uline{whose components satisfy a mild extra assumption besides being projective}.

Here I am struck by a slight discrepancy in terminology, as we
would rather have a correspondance
\diagram{
\[
\left\{
\begin{array}{l}
\text{integrators for B} \longrightarrow \text{abelianizators for A} \\[1ex]
\text{quasi-spe\xout{a}cial int\hand{$^{s}$}. for B} \longrightarrow \text{quasi-special abelian\hand{$^{s}$}. for A ,}
\end{array}
\right.
\]
}
and the same for \xout{x}"special" integrators and abelianizators. As I still
feel that the general \xout{\unsure{name of}} appelation of an \hand{``}integrator\hand{''}for \uline{any} projec-%
tive resolution of $\underline{Z}_B$ is adequate (without insisting that the components
should be in Addinf(B)), this kind of forces upon us to extend according-%
ly \inserted{still} the notion of an abelianizator for A. This does make a sense, using
the pairing (6) (which we had dismissed as an "irrelevant nicety" for
the time being" !) , and the corresponding equivalence

\margin{62}
